% Einheit 1 von 6 – Funktionswert und Punkt (bank/funktionsklassen-und-eigenschaften/e1.jsonl, zusammenbau v0.1)
%% TODO Zeitmarke Einheit 1: Marken-Zeile nicht lesbar
%% TODO Zweigzeile Teil 1: was hier gelernt wird (Satz in Schülersprache aus dem Einheitstitel) – Einheit 1
\einheitenkopf[e1]{Einheit 1 von 6 · Funktionswert und Punkt}
\zweigzeile{Abitur GK}
\setcounter{aufgabe}{8}

%% TODO Titel: Ich-kann-Satz fehlt in der Bank (Platzhalter: Kettenname) – Nr. 9
%% TODO Anweisung: ein Satz über der ersten Teilaufgabe fehlt in der Bank – Nr. 9
\begin{aufgabe}{Funktionswert und Punkt}
%% TODO \rechenplatz{n}: Schrittzahl des Grundfalls fehlt in der Bank (nur bei fester Schreibform, 2.3 b)
\begin{teile}
\teil Wie hoch ist der Ball $2$ s nach dem Abwurf? – Wert oder Stelle? \\ \kreuz{Stelle gegeben, Wert gesucht} \\ \kreuz{Wert gegeben, Stelle gesucht}
\teil $f(x) = x^3 - 3x^2 + 2$, Stelle $2$ – Funktionswert und Punkt? f(2) = \leerfeld , P( \leerfeld | \leerfeld )
\teil $f(x) = x^3 - 2x^2 - 5$ – liegt $P(3 | 4)$ auf dem Graphen?
\teil $f(x) = x^3 + 2x^2 - 4x + 6$ – Schnittpunkt mit der y-Achse und Steigung dort? S( \leerfeld | \leerfeld ), f'(0) = \leerfeld
\teil Tee kühlt ab: $T(t) = 20 + 70 \cdot e^{-0{,}1t}$ ($T$ in °C, $t$ in Minuten nach dem Aufgießen). Anfangstemperatur und Temperatur nach $10$ Minuten (auf eine Stelle)?
\teil Lies am Graphen ab: An welchen Stellen hat $f$ den Wert $3$?

\begin{ksys}[xmin=-4,xmax=4,ymin=-2,ymax=6,ablesen]\funktion{0.5*\x^2-1.5}{f}\end{ksys}
\teil $f(x) = x^2 + 1$ und $g(x) = 2x - 3$ – senkrechter Abstand der Graphenpunkte an der Stelle $3$?
\teil Modell $T(t) = 18 + 2t$ ($T$ in °C, $t$ in Stunden ab 9 Uhr); gemessen um 12 Uhr: $25$ °C. Prozentuale Abweichung des Modellwerts vom Messwert?
\teil $g(x) = -0{,}5x + 3$ – liegt $A(4 | 1)$ auf der Geraden? Zeichne die Gerade und $A$ ein.

\begin{ksys}[xmin=-2,xmax=6,ymin=-2,ymax=5]\end{ksys}
\teil $f(x) = -x^3 + 2x^2 + 2x + 3$. Liegt $P(-2 | 13)$ auf dem Graphen? Zeige, dass $3$ eine Nullstelle ist und dass $f$ keine weitere hat \hfill (FHR 2025).
\end{teile}
\end{aufgabe}

%% TODO Titel: Ich-kann-Satz fehlt in der Bank (Platzhalter: Kettenname) – Nr. 10
%% TODO Anweisung: ein Satz über der ersten Teilaufgabe fehlt in der Bank – Nr. 10
\begin{aufgabe}{Nullstellen und Werte: Werte a mit vorgegebener Lösungsanzahl von f(x) = a über die lokalen Extremwerte am Graphen angeben}
\begin{teile}
\teil Lies am Graphen ab: Für welche Werte $a$ hat $f(x) = a$ genau drei Lösungen?

\begin{ksys}[xmin=-5,xmax=5,ymin=-20,ymax=20,ystep=4,ablesen]\funktion{\x^3-12*\x}{f}\end{ksys}
\end{teile}
\end{aufgabe}

%% TODO Titel: Ich-kann-Satz fehlt in der Bank (Platzhalter: Kettenname) – Nr. 11
%% TODO Anweisung: ein Satz über der ersten Teilaufgabe fehlt in der Bank – Nr. 11
\begin{aufgabe}{Funktionswert und Punkt – Fehler finden, Begründen, Anwendung, Darstellungswechsel}
\begin{teile}
\teil Ich finde den Fehler: Gesucht ist die Stelle, an der $f(x) = 2x + 4$ den Wert $10$ hat. \rechnung{f(10) &= 2 \cdot 10 + 4 = 24}
\teil Begründe: Liegt $P(a | b)$ auf dem Graphen von $f$, dann gilt $f(a) = b$.
\teil Ein Handwerker berechnet seine Kosten mit $K(x) = 45x + 30$ ($K$ in €, $x$ Arbeitsstunden). Ein Kunde hat $250$ € eingeplant. Reicht das für $5$ Stunden?
\teil Der Graph zeigt $f(x) = -x^2 + 4x$. Lies $f(1)$ ab und bestätige den Wert am Term.

\begin{ksys}[xmin=-1,xmax=5,ymin=-2,ymax=5,ablesen]\funktion{-\x^2+4*\x}{f}\end{ksys}
\end{teile}
\end{aufgabe}
